\exerciseeditionsection{Solutions to Wald's Problems}
\ExerciseEditorialNote

\begin{enumerate}[leftmargin=2.8em,itemsep=1.2em]
\WaldSolution{1}
In the spherical, flat, and hyperbolic cases set, respectively,
\[
 r=\sin\psi,\qquad r=\psi,\qquad r=\sinh\psi.
\]
Then \(\dd\psi^2=\dd r^2/(1-kr^2)\), while the angular coefficient is
\(r^2\), giving the stated metric.  For \(k=+1\), \(r=\sin\psi\) fails to
distinguish \(\psi\) from \(\pi-\psi\).  A single such chart covers one
hemisphere of the spatial 3-sphere, with the equatorial 2-sphere at \(r=1\)
as a coordinate boundary.

\WaldSolution{2}
For the Robertson--Walker metric the nonzero curvature combinations reduce to
\[
 G_{ab}u^au^b=3\left(\frac{\dot a^2}{a^2}+\frac{k}{a^2}\right),
\]
and, on the spatial hypersurfaces,
\[
 G_{ab}h^a{}_ch^b{}_d
 =-\left(2\frac{\ddot a}{a}+\frac{\dot a^2}{a^2}
 +\frac{k}{a^2}\right)h_{cd}.
\]
With \(T_{ab}=\rho u_au_b+Ph_{ab}\), the temporal equation gives
\[
 3\frac{\dot a^2}{a^2}=8\pi\rho-\frac{3k}{a^2}.
\]
Combining the spatial equation with this result gives
\[
 3\frac{\ddot a}{a}=-4\pi(\rho+3P).
\]
The derivation depends on the spatial curvature only through
\({}^{(3)}R=6k/a^2\), so it covers both \(k=+1\) and \(k=-1\).

\WaldSolution{3}
\begin{enumerate}[label=\alph*),leftmargin=2em]
\item The modified Friedmann equations are
\[
 3H^2=8\pi\rho+\Lambda-\frac{3k}{a^2},
 \qquad
 3\frac{\ddot a}{a}=-4\pi(\rho+3P)+\Lambda.
\]
For a static solution,
\[
 \Lambda=4\pi(\rho+3P),
 \qquad
 \frac{k}{a^2}=4\pi(\rho+P).
\]
For ordinary matter this requires \(k=+1\) and \(\Lambda>0\).  For dust,
\[
 \Lambda=4\pi\rho=\frac1{a^2}.
\]
Conservation gives \(\rho a^3=\text{constant}\).  Put
\(a=a_0(1+\epsilon)\); then
\(\rho=\rho_0(1-3\epsilon)+\order(\epsilon^2)\), and the acceleration
equation becomes
\[
 \ddot\epsilon=4\pi\rho_0\epsilon=\frac{\epsilon}{a_0^2}.
\]
The exponentially growing solution shows the instability.

\item For \(T_{ab}=0\) and \(k=0\),
\(H^2=\Lambda/3\) and \(\ddot a/a=\Lambda/3\).  Thus
\[
 a(\tau)=a_0\exp\!\left(\pm\sqrt{\frac{\Lambda}{3}}\,\tau\right),
\]
the expanding or contracting flat slicing of de Sitter spacetime.
\end{enumerate}

\WaldSolution{4}
\begin{enumerate}[label=\alph*),leftmargin=2em]
\item Since \(u_a=-\nabla_a\tau\) and
\(g_{ab}=-u_au_b+h_{ab}\), direct evaluation of the Robertson--Walker
connection gives
\[
 \nabla_au_b=\frac{\dot a}{a}h_{ab}.
\]

\item For \(\omega=-u_ak^a\) and an affinely parametrized null tangent
\(k^a\),
\[
 \frac{\dd\omega}{\dd\lambda}
 =-k^ak^b\nabla_au_b
 =-\frac{\dot a}{a}h_{ab}k^ak^b
 =-\frac{\dot a}{a}\omega^2.
\]

\item Along the ray, \(\dd\tau/\dd\lambda=\omega\).  Dividing the preceding
equation by this relation gives
\(\dd\ln\omega/\dd\tau=-\dd\ln a/\dd\tau\).  Hence
\(a\omega\) is constant and
\[
 \frac{\omega_2}{\omega_1}=\frac{a(\tau_1)}{a(\tau_2)},
\]
which is Eq.~\eqref{eq:5.3.6}.
\end{enumerate}

\WaldSolution{5}
\begin{enumerate}[label=\alph*),leftmargin=2em]
\item Setting \(\dd s^2=0\) and suppressing the angular increments gives
\(\dd\tau=\pm a(\tau)\dd\psi\).  Therefore
\[
 |\Delta\psi|=\int_{\tau_1}^{\tau_2}\frac{\dd\tau}{a(\tau)}.
\]

\item For the closed dust model use conformal time \(\eta\):
\[
 a=A(1-\cos\eta),\qquad
 \tau=A(\eta-\sin\eta),\qquad 0\leq\eta\leq2\pi.
\]
Since \(\dd\tau/a=\dd\eta\), the ray travels
\(\Delta\psi=2\pi\), one complete great circle.

\item For the closed radiation model,
\[
 a=A\sin\eta,\qquad
 \tau=A(1-\cos\eta),\qquad 0\leq\eta\leq\pi.
\]
Again \(\dd\tau/a=\dd\eta\), so \(\Delta\psi=\pi\), half a great circle.
\end{enumerate}
\end{enumerate}
